\section{Q1d: Test statistics}
\label{sec:Q1d}
The output of the test statistics is already shown in the previous two sections. This section serves as an explanation, and highlight of the code involved.\\
\\
To assess the goodness of our fits, we perform a G-test, given by
\begin{equation}
    G = 2\cdot\sum_i^{N_\text{bins}-1}O_i\log\left(\frac{O_i}{E_i}\right),
\end{equation}
where $O_i$ is the integer observed count, and $E_i$ is the expected count. Since the observed and expected model should integrate to the same value (the total number of observations), we rescale our expected model $E_i$, such that this is satisfied.\\
\\
We calculate the significance $Q$ using $Q=1-P(x, k)$, where $P(x,k)$ is the $\chi^2$ probability density function, given by
\begin{equation}
    P(x,k) = \frac{\gamma\left(\frac{k}{2}, \frac{x}{2}\right)}{\Gamma\left(\frac{k}{2}, \frac{x}{2}\right)}.
\end{equation}
Here, $\gamma(k, x)$ is the incomplete lower gamma function, $\Gamma(k, x)$ is the gamma function, $x$ is the test statistic ($G$ in our case), and $k$ is the number of degrees of freedom. Naively, $k$ is given by $N_\text{bins} - 3$, as we have 3 (independent) fitting parameters. However, not all bins are independent; once $N_\text{bins}-1$ bins have been filled, the number of points in the last bin is fixed. As such, the total number of degrees of freedom is given by $k=N_\text{bins}-1-3 = N_\text{bins}-4$.

The code for the G-test and the $\chi^2$ distribution if given in \verb|helperscripts/stat.py|:
\lstinputlisting[style=pystyle, language=Python, firstline=40]{../helperscripts/stat.py}
